\section{从绍兴到崖山：战争怎样改变道路、财政与选择}

\subsection{南渡战争中的河道和军队}

建炎年间的战场并没有一条稳定前线。金军的行动、宋廷的驻跸、勤王军的来去和地方的自保，都围绕渡口、城镇与水道展开。金军渡江而后北撤，\eventref{SS-1130-EVENT-114}{并不证明长江从此成为不可跨越的天险}；它说明在那一段时间内，船只、补给、抵抗和更大的战略考虑使继续推进成本上升。对于南宋朝廷，长江的价值也不只在于阻挡敌军，它还是官员、粮食、军器和流民必须使用的通道。

军队整编的困难由此出现。能够在某一次会战取胜的将领，不必然能够长期维持与中枢、地方、粮道和其他部队的关系；朝廷担心的也不只是敌军，还包括兵权、财政和内部秩序怎样共存。苗刘兵变让行在看见禁军与宫廷关系的危险，岳飞等部队的活动则让恢复战场与中枢控制的矛盾更尖锐。后世常把这些事件铺成忠奸对决，然而《要录》所见的任免、奏议和军政处置提醒我们，赏给、驻地、补给与朝廷对战局的不同判断同样在塑造结果。

一一四〇年前后的反攻和次年的和议，不应以“本可成功”或“必然失败”倒叙。宋方与金方战报对军数、战果和责任的夸张并不相同，后来的纪传又常把复杂的组织问题压到少数人物身上。可以较稳妥地说，南宋需要在江淮、襄汉与中枢之间协调军粮、赏给和命令，金也要判断它在南方推进的承受力；但具体兵额、每一次战斗的伤亡和私人动机，不能越过材料的界限写得确定。

\subsection{和议是一种持续的军事安排}

绍兴和议\eventref{SS-1141-EVENT-120}{给南宋带来的不是解除武装的和平}。它设置了使节、岁币和边界的关系，也让国家能够把一部分注意力转回官制、财赋和城市。然而两淮、川陕和江防的守备并没有消失。边境上的堡寨、渡口和市集既可能成为贸易和接触点，也可能在紧张时迅速变成侦察、封锁和征发的节点。

和约文本的仪式语言必须同它的执行分开。称臣、兄弟、赐予或誓书说明谈判双方以何种词汇处理关系，却不能凭自身告诉我们哪一座城中的税收、司法或居民认同已经改变。宋方材料和金方材料可彼此校验使节往返、军事冲突和大致年序；对路线、兵数、贡赋和盟辞的细节，仍须保留传本与立场差异。将外交写作一方的屈服或另一方的征服，都会把沿边社会日常承受的守备、贸易与不确定性排除。

一一六一年南下的金军再次测试长江防线。采石等地的战事之所以重要，在于它把河面、船只、岸防和补给集中到同一场行动里。战后高宗禅位，孝宗即位，朝廷面临的是既要利用战局变化、又不能忽略兵粮与组织承受力的难题。符离北伐的受挫和隆兴和议的达成，显示决策并非在“恢复”与“屈服”之间抽象选择，而是要在情报不全、部队协同有限、河道运输和对方反应都不确定的条件下决定风险由谁承担。

\subsection{开禧北伐：动员不能由口号替代}

宁宗朝的开禧北伐\eventref{SS-1206-EVENT-132}{把长期积累的财政和军事问题重新推向前线}。准备远征时，朝廷要依靠地方的粮运、兵员、情报和道路，金方也会调整边防与外交。行动未能按预想推进，韩侂胄被杀，嘉定和议随之结束战事。把失败只归于一位权臣，容易省略地方军队、对金情报、后勤能力和中枢派系如何相互限制；把它只归于财政不足，又会忽略战场和政治判断并不是由一张收支表自动决定。

这次北伐特别能说明“军费”不是一个孤立项目。为军队筹集的米、钱、布帛和船只，要由州县征集、储存并运到不同前线；征集造成的压力又会影响市场、债务和地方合作。制度文献可以追到哪些名目被提出，不能替我们量化每一处家庭的负担。地方社会也不是静止地承受：有人承运、有人供给、有人规避、有人在战争中失去原有的生计。只是后几类行动留在材料中的比例，很难与官员和将领相同。

嘉定和议后边境并未归于无事。金军对四川、两淮的进攻，以及南宋持续的驻防投入，都使“和平”成为需要不断维持的状态。边区人们的选择不一定与临安的论争一一对应：一座城要考虑粮仓和守军，一处渡口要考虑通行和税收，家庭要考虑田地、服役和逃难的可能。让这些不同尺度并存，才能避免把战争史只写成朝廷的决策史。

\subsection{攻金协作后的空白地带}

一二三四年蔡州陷落后，南宋曾试图进入河南故地，\eventref{SS-1234-EVENT-138}{短暂的宋蒙协作随即暴露出彼此不同的目标}。共同攻击金并没有产生共同管理道路、城池和粮区的方案。对南宋而言，恢复旧地既有政治象征，也需面对久经战乱地区的粮食、居民和防务；对蒙古而言，军事推进形成的控制关系未必接受宋廷的规划。将这一转折写作“联蒙灭金”的一句话，会把协作结束时的空间竞争隐藏起来。

蒙古开始攻宋后，战争并非从北向南整齐推进。四川、荆襄、两淮和沿江各有不同的水陆条件、城防体系和地方政治。蒙古军的机动能力、攻城技术与区域盟友构成压力，南宋的山城、江防、舰船和地方供给也形成抵抗。宋、元两边的叙事尤其容易把对方写成单一行动者；实际的行军、投降、守城、转运和谈判，发生在许多并不同时的地点。

一二五九年钓鱼城战事和蒙哥去世改变了战争节奏。关于蒙哥死因的说法来源复杂、互有冲突，不能以传奇细节填补材料空白。较可靠的观察是，重庆附近的山城防御和峡江交通使蒙古对宋的战争不能按预设节奏轻易完成。城址与地形能够使我们看清防御为何困难，却不能把高处、城墙或某件兵器直接变成胜败的唯一原因。防守的持续仍依赖粮、水、人力和与周边区域的联系。

\subsection{襄樊：围城的时间比一场决战长}

襄阳、樊城从一二六八年起受到长期围困，\eventref{SS-1268-EVENT-144}{汉水和江陵方向的交通成为双方争夺的关节}。围城并不是把军队摆在城外等待城门开启；封锁者必须维持营地、船队和补给，守方则须让粮食、援军和消息穿过水陆网络。汉水连到江汉平原，江陵及长江又把这一地区同江南财赋连接。任何一段被切断，城中的存粮、兵员和选择就会改变。

宋末叙事常突出守城忠义，元方叙事则强调攻取与统一；两者都保存了重要的行动线索，也都服务各自的政治记忆。长围期间的军数、伤亡和城内生活尤其不宜依单方数字断言。我们能够看见的是，围城拖延多年而非瞬间胜负，因而给周边村镇、江路商运和地方征发造成持续压力。没有留下记录的人如何调整生活，不能凭忠义叙事或胜利叙事替他们决定。

一二七三年襄樊失守\eventref{SS-1273-EVENT-145}{使长江中游的防御结构受到难以补上的破口}。这并不表示南方所有城市立即失去抵抗能力，而是攻击者获得了更有利的水陆通道，南宋则失去一个长期吸收冲击的节点。随后元军南下、丁家洲失利、临安局势恶化，说明此前依赖的运输与政治调度已经难以恢复。把这一连串事件写成守城者道德高低的结论，反而会遮蔽城、河、船和粮道共同决定的结构性压力。

\subsection{临安降元与南方的不同时间}

临安降元\eventref{SS-1276-EVENT-148}{首先是宋廷中枢的政治转换}。皇帝、太后、官员和文书被置于新的安排中，并不意味着东南、福建、岭南和海上的军事与社会关系在同一天改变。宗室和部分军民继续南行，地方有人守、有人降、有人观望，更多人仍须处理粮食、亲属和生计。宋末“遗民”与元初“归附”的称谓是后人认识这段历史的重要入口，但它们也会把复杂、反复而地方化的选择压成固定身份。

元军的接收同样不是抽象的行政动作。新的官府需要税源、道路、船只和地方中介；愿意或不得不合作的人可能以不同条件进入新秩序。后来形成的元朝档案和史书能记录一部分任命、赋役和军政安排，却有胜利者的制度视角。它们不能回溯为所有地方在一二七六年已经完成转换。与宋方材料并读的价值，正在于让朝廷终结、地区控制和居民生活的三种时间不被混为一刻。

\subsection{崖山的海面与战争的余波}

海上的宋军和宗室继续行动，最终在一二七九年崖山覆亡。\eventref{SS-1279-EVENT-150}{这场终局需要放回珠江口与沿海补给的空间}。舰船、淡水、粮食、潮汐、岛屿和人员构成了海上政权能够维持多久的条件。宋、元叙事留下了强烈的忠义与胜利语言，后来的文学和纪念又使崖山成为象征。象征并非无价值，却不能让读者误以为所有船上和沿岸的人共享同一种选择或同一种记忆。

海战也不应被解释为技术优劣的一次展示。船只需要木材、工匠和港口，水手与军人需要食物、疾病防护和指挥，舰队能否移动还受季节和地形限制。现存文本很难精确复原每一艘船的位置、全部人数和每位行动者的动机；考古和海洋地理可以使某些条件更具体，却不能补出一幅无缝的战场图。把不确定性说清楚，反而能看见战争为何对沿海社会留下较长的余波。

南宋的军事史至此不是一段由英雄或失败者单独书写的尾声。从行在的渡口到两淮的寨堡，从四川山城到襄樊水道，再到珠江口的舰船，战争反复要求国家把人、粮、钱、船和文书接起来。它有能力在很长时间里完成这种连接，也因连接集中在有限的水路、地区和中介身上而不断付出代价。这一代价在财政、家庭、港口和寺院材料中留下碎片；把碎片接回年序，才能理解政权的终结并不取消社会继续面对的道路与资源问题。
